\section{The irreducible code on 176 residual blocks}
Here $H=H_1$, $C=C_{176}=H^{\mathsf T}(e^\perp)$ and $\dim C=21$ by (1).
The integral Gram matrix is $40I_{22}+16J_{22}$, reducing to $I_{22}+J_{22}$.
The common incidence theorem proves LCD and the two complement--dual
identities. The constant vector belongs to $C^\perp$, because its image
under $H$ is $56e$, which annihilates $e^\perp$. Thus the full orthogonal
decomposition in the source is established before any module classification.

\subsection{The exact module and group certificate}
The literal permutations in \texttt{module\_data.json} preserve every
adjacency entry of $\Gamma$. Their closure, computed by exact permutation
composition, has order $443520$. They also permute the $22$ point-incidence
rows of $H$. The resulting faithful degree-$22$ action is explicitly
identified with \texttt{MathieuGroup(22)} by an isomorphism checked in GAP.
Thus the subgroup $G$ they generate is $M_{22}$. This conventional group
identification is distinct from proving that $G$ is the \emph{full} code
automorphism group. The latter will follow from the common refinement theorem \cite[Theorem~6.4]{U}.
The auxiliary group calculation also verifies that the degree-$176$ action
is primitive with stabilizer $A_7$ of order $2520$.

The row-reduced adjacency basis is checked byte for byte against the basis
of the retained complete enumeration. For each coordinate generator, its
images on these $21$ rows give a $21\times21$ representation matrix.
The checker reconstructs those images from the proposed matrices, so the
matrix module is the actual code $C$, not a representation with only the
same dimension. Deterministic multiplication and elimination produce
$441$ independent generator words, all of length at most eight. Each word
is evaluated again and the $441$ flattened matrices are row-reduced.
Their rank is $441$. A separate GAP calculation confirms the same generated
algebra dimension. The common matrix-algebra theorem \cite[Theorem~6.2]{U} proves that $C$ is absolutely
irreducible. This action is faithful as well: the $176$ columns of a code
basis are distinct, so a coordinate permutation fixing every codeword
fixes every coordinate. Column distinctness is checked in the certificate.

The retained complete enumeration gives $d(C)=56$ and exactly $924$ words
of weight $56$. The supplied literal weight-$56$ word belongs to $C$.
Its orbit under $G$ has $462$ words and is disjoint from the orbit of its
negative. Together those two orbits have $924$ words, so they exhaust the
minimum words. Their supports are the same $462$ distinct subsets.
These assertions are checked by finite orbits and the already accepted
count; no new exhaustive codeword computation is involved.

For this support-incidence graph, with $176+462$ vertices, the deterministic
refinement rule of the common refinement theorem \cite[Theorem~6.4]{U} successively individualizes
coordinates $1,5,29,61,2$. The corresponding cell sizes are
\[
 176,\quad70,\quad6,\quad2,\quad3.
\]
All coordinate colours then differ. Therefore
\[
 |\Aut(C)|\leq176\cdot70\cdot6\cdot2\cdot3=443520.
\]
Every graph automorphism preserves its adjacency row space, so
$G\leq\Aut(\Gamma)\leq\Aut(C)$. The already checked order of $G$ forces
equality throughout:
\[
 \Aut(C)=\Aut(\Gamma)=G\cong M_{22}.                 \tag{10}
\]
This upper bound uses no graph-isomorphism backend. Such software was
used only to discover the two literal generators.

The complete retained weight distribution determines the dual by the common MacWilliams theorem \cite[Theorem~7.1]{U}.
The certificate computes all coefficients exactly, checks nonnegative
integrality and total $3^{155}$, and obtains zero at weights one through
five and $184800$ at weight six. Thus the parameters are
$[176,21,56]_3$ and $[176,155,6]_3$, with the asserted minimum counts.

\subsection{The precise representation boundary}
Theorem~1(c) describes the \emph{unique $21$-dimensional} irreducible
$F_3$-module for $M_{22}$. Proposition~4's final sentence omits the dimension,
but its proof immediately restores it and invokes the known $3$-modular
character table. We retain that intended classification statement, and
the proof's further assertion that $21$ is the smallest nontrivial
irreducible dimension, from the source's reference [13]:
Jansen--Lux--Parker--Wilson, \emph{An Atlas of Brauer Characters} (1995).
These assertions are understood up to module isomorphism. Our full
matrix-algebra calculation proves irreducibility of this particular
module; it does not classify all modules. The unqualified literal
uniqueness sentence cannot be correct, since the trivial one-dimensional
module is irreducible too.

The source's historical near-optimality comparison for the dual is a
separate assertion based on its tables, not a current table claim or a
consequence of (10). No additional irreducibility is asserted for the
$155$-dimensional dual, which contains the constant line. The group in
(10) consists of coordinate permutations; allowing scalar monomial maps
would change the automorphism convention.

